\documentclass[10pt,a4paper]{amsart}
\usepackage[margin=19mm]{geometry}
\usepackage{amsmath,amssymb,mathtools}
\usepackage[scr=rsfs,cal=euler]{mathalfa}
\usepackage{xfrac}
\usepackage[unicode,hidelinks,bookmarks=false]{hyperref}
\newcommand{\ie}{i.e.\ }
\newcommand{\cf}{cf.\ }
\newcommand{\resp}{respectively\ }
\newcommand{\Subsection}{\subsection}
\providecommand{\nobreakdash}{\nobreak\hbox{-}}
\setlength{\emergencystretch}{2em}
\multlinegap0pt
\usepackage{bm}
\usepackage{bbm}
\usepackage{dsfont}
\usepackage{xspace}
\usepackage{cancel}
\usepackage{mathtools}
\usepackage{amsthm}
\makeatletter
\@ifundefined{theorem}{\newtheorem{theorem}{Theorem}}{}
\@ifundefined{proposition}{\newtheorem{proposition}[theorem]{Proposition}}{}
\makeatother
\def\r{\mathbb R}
\def\v{\vec{v}}
\title[A Note on Helicoidal Singular Minimal Surfaces]{A Note on Helicoidal Singular Minimal Surfaces}
\author{Rafael L\'opez}
\date{Source: arXiv:2507.13669v1 (CC BY 4.0)}
\begin{document}
\maketitle

\noindent\emph{Source and licence.} A Note on Helicoidal Singular Minimal Surfaces. Version arXiv:2507.13669v1 (CC BY 4.0), distributed under the Creative Commons Attribution 4.0 International licence, as recorded in the arXiv licence metadata for this exact version and shown on the abstract page https://arxiv.org/abs/2507.13669v1. Full rights notice: licenses/arxiv-2507.13669-CC-BY-4.0.txt.\par\smallskip

\noindent\textit{Editorial scope.} Counted unit: the Proposition whose source label is \texttt{pr1} in the article (source0.tex lines 164--166, source label \texttt{pr1}), delivered together with the article's own complete original proof of it (source0.tex lines 168--190, one \texttt{proof} environment of 23 lines). No mathematics is written, repaired or re-derived: statement and proof are the article's own text line for line. Required context retained uncounted: eq1 (lines 76--78); para (lines 116--118); arc (lines 142--147); mean2 (lines 149--151); normal2 (lines 153--155); reg (lines 158--160). Every definition, display and label of those lines is retained unchanged. Omitted from this package and not redistributed: 1--75, 79--115, 119--141, 148--148, 152--152, 156--157, 161--163, 167--167, 191--317. Nothing inside a retained excerpt is omitted or altered: every retained body is delivered exactly as the article's lines are written, so no omission receipt of any kind accompanies this package. Citations: the retained excerpts carry no citation command at all, so no bibliography entry of the article is transcribed and no reference is redistributed. Reference adaptation: none was needed and none was made; every reference inside the delivered unit resolves inside it, so no unresolved reference or citation is produced. Printed slips kept literal: no printed word, symbol, hypothesis or number of the counted statement or of its proof was altered. Layout and class: the article's own class is \texttt{amsart} with packages including amssymb, amsthm, amscd, graphicx, color, url; the wrapper is the lane's pinned \texttt{amsart} editorial wrapper at 10pt on A4, which restates the theorem-like environments the retained text needs and adds the article's own macro definitions that the retained text needs (\texttt{\textbackslash r}, \texttt{\textbackslash v}), with extra wrapper packages (bm, bbm, dsfont, xspace, cancel, mathtools). The delivered numbering is the unit's own. Assets: no retained span contains a figure environment, a graphics-inclusion command, a PDF-page inclusion, a table or a TikZ picture; no image file is copied into this package, the package declares no asset dependency and no image file is redistributed with it. Licence: version arXiv:2507.13669v1 of this article is distributed under the Creative Commons Attribution 4.0 International licence, as recorded in the arXiv licence metadata for this exact version and shown on the abstract page https://arxiv.org/abs/2507.13669v1. A full rights notice is delivered at licenses/arxiv-2507.13669-CC-BY-4.0.txt. Characters: every retained excerpt is pure ASCII, so the delivered engine, XeLaTeX with its default Latin Modern text font, reports no missing character.
\begin{equation}\label{eq1}
H(p)=\alpha\frac{\langle N(p),\v\rangle}{\langle p,\vec{v}\rangle},\quad p\in\Sigma,
\end{equation}

\begin{equation}\label{para}
\Psi(s,t)=\mathcal{H}_t(\gamma(s))=(x(s)\cos t,x(s)\sin t,z(s)+ht).
\end{equation}

\begin{equation}\label{arc}
\begin{split}
x'(s)&=\cos\theta(s),\\
z'(s)&=\sin\theta(s).
\end{split}
\end{equation}

\begin{equation}\label{mean2}
H=  \frac{x (x^2+h^2) \theta '+\sin\theta  \left(x^2+2 h^2 \cos ^2\theta \right)}{ (x^2+h^2 \cos ^2\theta )^{3/2}}.
\end{equation}

\begin{equation}\label{normal2}
N=\frac{1}{\sqrt{x^2+h^2\cos ^2\theta}}\left(h\cos\theta\sin t-x\sin\theta\cos t,-h\cos\theta\cos t-x\sin\theta\sin t,x\cos\theta\right).
\end{equation}

\begin{equation}\label{reg}
EG-F^2=x^2+h^2 \cos ^2\theta>0.
\end{equation}

\begin{proposition} \label{pr1}
Let $\Sigma$ be a   surface invariant by the helicoidal group $\{\mathcal{H}_t\colon t\in\r\}$ with pitch $h$.  If $\Sigma$ is a singular minimal surface with $\vec{v}=(0,0,1)$, then $h=0$.
\end{proposition}

\begin{proof} From \eqref{para} and \eqref{normal2} we have 
\begin{equation*}
\begin{split}
\langle N,\vec{v}\rangle&=\frac{x\cos\theta}{\sqrt{x^2+h^2\cos^2\theta}},\\
\langle p,\vec{v}\rangle&=\langle \Psi(s,t),\vec{v}\rangle=htz.
\end{split}
\end{equation*}
By using the expression of the mean curvature $H$ in \eqref{mean2}, equation \eqref{eq1} writes as a polynomial on $t$ of degree $1$, namely, 
$$A_0(s)+A_1(s) t=0.$$
Then the coefficients $A_0$ and $A_1$ must be zero. A computation of these coefficients gives
\begin{equation*}
\begin{split}
A_0&=z \left(x (x^2+h^2) \theta '+\sin\theta (x^2+2 h^2 \cos ^2\theta\right))- \alpha  x \cos \theta \left(x^2+h^2 \cos ^2\theta\right),\\
A_1&=h \left(x (x^2+h^2) \theta '+\sin\theta (x^2+2 h^2 \cos ^2\theta )\right).
\end{split}
\end{equation*}
The linear combination $h A_0-z A_1=0$ becomes
$$\alpha h x\cos\theta(x^2+h^2\cos^2\theta)=0.$$
Since $x^2+h^2\cos^2\theta\not=0$ by \eqref{reg}, we deduce $h=0$, or $x=0$ identically or $\cos\theta=0$ identically. If $h=0$, then the result is proved. Assume now that $h\not=0$ and we see that the cases $x=0$ or $\cos\theta=0$ identically cannot occur. 

If $x=0$ identically, then $x'=\cos\theta=0$ hence $x^2+h^2\cos^2\theta=0$, which it is not possible. If $x'=\cos\theta=0$ identically, then $x=x(s)$ is a constant function, $x(s)=x_0$ with the extra condition that $x_0\not=0$ by  regularity \eqref{reg}. In particular, from \eqref{arc} we deduce $z(s)=\pm s+z_0$. Computing again  the coefficients  $A_0$ and $A_1$, we simply have 
$A_0(s)=x_0(\pm s+z_0)$ for some $z_0\in\r$ and $A_1(s)=hx_0$. Then $A_0=0$ gives a contradiction. This completes the proof of Prop. \ref{pr1}.
\end{proof}

\end{document}
